% 本文件由 LaTeX 排版 Agent 根据有效 translation.json 自动生成。
% author=Councils; work_id=3818; title=迦太基公会议（公元257年）

\chapter{迦太基公会议（公元257年）}

% structure: p0001 source=h1 target=omitted reason=page-title-already-rendered-by-parent

\Emph{在迦太基举行的主教会议，由伟大的圣殉道者、迦太基主教西彼廉主持。}\par

当许多主教于九月朔日在迦太基聚集——他们来自阿非利加省、努米底亚和毛里塔尼亚，并有司铎与执事陪同（大部分民众也同样在场）——当犹巴伊安写给西彼廉的圣函、西彼廉就异端者的圣洗问题给犹巴伊安的覆函，以及同一位犹巴伊安随后写给西彼廉的信函被宣读之后，\par

西彼廉说：\Quote{你们已经听见了，我亲爱的同僚们，我们的同僚主教犹巴伊安如何征询我卑微的意见，涉及异端者非法且亵圣的圣洗，以及我给他的答复；一如我们先前所持的见解，即来到教会的异端者应受教会的圣洗，并藉此得以圣化。此外，还向你们宣读了犹巴伊安的另外一封信，他在信中以其真诚而虔诚的敬慕回应我们的信函，不仅表示赞同，还感恩自己因此受到了教诲。因此，现在只待我们每一位逐一说出对此事的见解，不谴责任何人，也不将与我们意见不同的人罢黜共融之权。}\par

\Quote{因为没有一个人将自己立为（众主教的）主教，或试图以专制的恐惧强迫同僚服从他，因为每位主教，根据自由与权力的许可，都拥有自己的意愿，他既不能被他人审判，也不能审判他人。但我们等待我们普世的主、我们的主耶稣基督的审判，祂是惟一且独一有权柄者，既能提拔我们治理祂的教会，又能审判我们在那职位上的行为。}\par

\Emph{主教们随后逐一宣告反对异端者的圣洗。最后（col. 796）：}\par

迦太基的宣信者与殉道者西彼廉说：\Quote{那封写给同僚犹巴伊安的信函已充分陈述了我的见解：那些根据福音与宗徒的见证被称为基督的仇敌与假基督的异端者，当他们来到教会时，应以教会的\OriginalTerm{唯一}{unico}圣洗受洗，好使他们由仇敌变为朋友，由假基督变为基督徒。}\par

\SourceRecord{Translated by Henry Percival.；Nicene and Post-Nicene Fathers, Second Series；Vol. 14.；Edited by Philip Schaff and Henry Wace.；Buffalo, NY: Christian Literature Publishing Co.,；1900.；<http://www.newadvent.org/fathers/3818.htm>.}
